\section{Introduction}\label{sec:introduction}

A unital ring $R$ is \emph{directly finite} if $ab=1$ implies $ba=1$ for all
$a,b\in R$. For a field $K$ and a group $G$, the group algebra $K[G]$ consists
of finite formal sums of elements of $G$ with coefficients in $K$, with
multiplication extending the group operation bilinearly. Kaplansky's direct
finiteness conjecture asserts that every such group algebra is directly finite.
We prove that the implication can fail even when the field is finite and the
group is finitely presented and torsion-free.
A classifying complex for $G$ is a connected CW complex with fundamental group
$G$ and contractible universal cover.

\begin{theorem}\label{thm:main}
There exist a finitely presented torsion-free group $G$ and elements
$a,b,c\in\Ftwo[G]$ such that
\[
 ab=1,\qquad ac=0,\qquad c\ne0.
\]
Consequently $ba\ne1$. The group $G$ admits a finite two-dimensional
classifying complex.
\end{theorem}

The auxiliary element $c$ certifies failure of the reverse identity:
if $ba=1$, then $c=(ba)c=b(ac)=0$.

\subsection*{History and related work}
Kaplansky's formulation separates the characteristic-zero theorem from the
positive-characteristic question~\cite[Part~II, Section~3, Problem~2,
pp.~122--123]{Kaplansky1972}. In characteristic zero, group algebras are directly
finite for every group. The operator-algebra argument also applies to matrix
rings; Montgomery gave a concise $C^*$-algebra proof~\cite{Montgomery1969}.
A ring whose matrix rings of every finite size are directly finite is called
\emph{stably finite}.

In arbitrary characteristic, Ara, O'Meara, and Perera proved stable finiteness
for group algebras of free-by-amenable groups over division
rings~\cite{AraOMearaPerera2002}. Elek and Szab\'o extended the conclusion to
sofic groups~\cite[Corollary~4.7]{ElekSzabo2004}, a class defined by finite
approximate permutation actions. Their theorem implies that the group in
Theorem~\ref{thm:main} is not sofic. A complementary approach studies the
finite cancellation patterns forced by a one-sided inverse. Dykema, Heister,
and Juschenko organized these patterns into universal groups and established
direct finiteness for several small support pairs over
$\Ftwo$~\cite{DykemaHeisterJuschenko2015}.

Two companion manuscripts are especially relevant. A separate construction
gives a scalar counterexample over a finite field of characteristic two for
a finitely presented group with odd-order torsion~\cite[Theorem~1.1]{OpenAI2026DF}.
Theorem~\ref{thm:main} uses the prime field $\Ftwo$ and a torsion-free group.
Its method develops the graph and cone construction of
\emph{A Torsion-Free Group Algebra with Zero Divisors}~\cite{OpenAI2026ZD}.
That manuscript supplies the parity mechanism for a zero product and the
arrangement-exclusion strategy. We adapt those arguments to produce a
one-sided inverse and reproduce the estimates and topology needed here.

Graphical presentations relate coincidences of labeled paths to geometric
properties of quotient groups, as in Gromov's random graphical
constructions~\cite{Gromov2003} and Ollivier's small-cancellation
theorem~\cite{Ollivier2006}. Bereznyuk's spherical-diagram criterion for
graphical asphericity distinguishes coincident graph lifts from genuinely
different ones~\cite{Bereznyuk2019}. The same distinction is essential here:
pairings of two occurrences of the same graph edge are removed by cone
surgeries. We prove the required asphericity and root protection directly in
Section~\ref{sec:topology}.

\subsection*{The construction and its main estimates}
The group is obtained from a bouquet of circles by attaching cones on finite
labeled graphs. A cone kills every closed path in its graph, so a vertex in a
chosen component determines a group element by reading a path from a fixed
root to that vertex. The desired algebra elements are finite sums of these
path labels and their inverses.

Two graphs are used. A common outgoing label allows a simultaneous step in
both coordinates of their vertex products. The group element contributed to
the corresponding algebra product is constant along such a step. If every
vertex of a simultaneous-step component has odd degree, the component has
even cardinality and its contributions cancel over $\Ftwo$. A single
even-degree exception instead leaves one surviving contribution. We prescribe
the outgoing label sets so that the first graph paired with itself has exactly
this exception, while the first graph paired with the second has none.

The change to the parity construction in \cite{OpenAI2026ZD} is to produce
product $1$ as well as product $0$. Seven additional labels, organized by complements of
lines in the Fano plane, create the required exception while keeping every
inverse-label incidence count balanced. This changes the random model, and
we establish its counting estimates directly.

The main difficulty is to make the cone attachments aspherical while ensuring
that a path from a root to a different graph vertex has nontrivial image in
the group. Both conclusions follow by excluding certain planar pairings of
graph paths. The argument has two parts. First, a weighted estimate rules out
bounded collections of paths with a small proportion of unpaired positions.
Repeated traversal of a graph edge is the main counting obstacle: two
compared path words need not impose two sets of random edge conditions.
We organize the count by the number of times each edge is used, and choose
subdivisions that nearly agree under all the comparisons. This makes the
decay of word weights available despite repeated traversals. The permitted
unpaired proportion is fixed before the bounds on the number and length of
paths. Second, a planar separator extracts one of these bounded collections
from any spherical arrangement, however large. The order of these parameter
choices is what allows a local counting estimate to exclude all arrangements.

We give the graph criterion and the parity calculation in
Section~\ref{sec:algebra}. Section~\ref{sec:random} constructs the vertex types
and proves the word, girth, and diameter estimates.
Section~\ref{sec:patterns} proves the bounded-path estimate, and
Section~\ref{sec:planar} carries out the planar extraction.
Section~\ref{sec:topology} proves the topological criterion by cone surgeries.
Section~\ref{sec:assembly} then assembles these conclusions.
